\chapter{Stable-homotopy interfaces absent from Mathlib}
\label{chap:stable-background}

Mathlib v4.32.2 has preadditive and triangulated-category infrastructure, but
no model of spectra, the Steenrod algebra, the Adams spectral sequence, or
$H\mathbb F_2$-synthetic spectra.  Following the project boundary, these are
abstract interfaces with exactly the operations and laws used by the paper.

\section{Spectra, triangles, and homology}

% SOURCE: PROJECT_BOUNDARY.md, Confirmed design decision 1; aimpaper/main.tex:253-263, 670-758.
\begin{definition}[Stable homotopy context]
  \label{def:stable-context}
  \uses{def:mathlib-shift,def:mathlib-triangle,def:mathlib-pretriangulated,def:mathlib-monoidal-category,def:mathlib-symmetric-category}
  \notready
  A stable homotopy context consists of a pointed additive homotopy category
  of spectra, suspension equivalences $\Sigma^n$, a symmetric monoidal smash
  product with sphere unit $S^0$, cofiber sequences, and distinguished
  triangles.  It includes rotation, functoriality, smash compatibility, and
  the sign and suspension conventions needed below.  Stable homotopy groups
  and their exact sequences are exposed separately so that later proofs do not
  hide them inside this context record.
\end{definition}

% SOURCE: aimpaper/main.tex:126-138 and 253-263; suspension conventions used throughout.
\begin{definition}[Sphere spectrum and suspension]
  \label{def:stable-sphere-suspension}
  \uses{def:stable-context}
  \notready
  The monoidal unit is the sphere spectrum $S^0$.  For every $n\in\mathbb Z$,
  suspension gives an equivalence $\Sigma^n$ with coherent isomorphisms
  $\Sigma^0X\simeq X$ and $\Sigma^{m+n}X\simeq\Sigma^m\Sigma^nX$, compatible
  with smash products and mapping groups.
\end{definition}

% SOURCE: aimpaper/main.tex:102-170 and throughout Sections 2--7.
\begin{definition}[Stable homotopy groups]
  \label{def:stable-homotopy-group}
  \uses{def:stable-sphere-suspension}
  \notready
  For a spectrum $X$ and $n\in\mathbb Z$, define
  \[
    \pi_nX=[\Sigma^nS^0,X].
  \]
  These are abelian groups, functorial in $X$, with suspension and composition
  isomorphisms fixed by the context's sign convention.
\end{definition}

% SOURCE: aimpaper/main.tex:253-263, global standing assumptions.
\begin{definition}[Admissible spectrum]
  \label{def:admissible-spectrum}
  \uses{def:stable-context}
  \notready
  A spectrum is admissible when it is $2$-complete, connective, and of finite
  type.  Strong convergence is deliberately not part of this predicate:
  whenever a classical or synthetic Adams spectral sequence is invoked, its
  convergence witness is supplied separately.  This avoids defining
  admissibility in terms of spectral sequences that are themselves indexed by
  admissible spectra.
\end{definition}

% SOURCE: aimpaper/main.tex:135-159, where products are used for the sphere;
% standard associative-algebra-object structure in a symmetric monoidal category.
\begin{definition}[Unital multiplication data]
  \label{def:unital-multiplication-data}
  \uses{def:stable-context,def:stable-sphere-suspension}
  \notready
  A unital multiplication datum on a spectrum $R$ is an associative algebra
  object in the stable symmetric monoidal context: it includes multiplication
  $R\wedge R\to R$, a unit $S^0\to R$, and all associativity and unit
  coherences.  This is extra structure and is not implied by admissibility.
  The sphere $S^0$, as the monoidal unit, has its canonical such datum.  The
  same terminology will be used for associative algebra objects in the
  synthetic context.
\end{definition}

% SOURCE: aimpaper/main.tex:688-769; Pstragowski interface used in Proposition 3.1.
\begin{definition}[$H\mathbb F_2$ homology]
  \label{def:hf2-homology}
  \uses{def:stable-context,def:mathlib-f2}
  \notready
  The context contains an Eilenberg--Mac Lane ring spectrum $H\mathbb F_2$ and
  a homology functor $H_*(-;\mathbb F_2)$ from spectra to graded
  $\mathbb F_2$-vector spaces.  It preserves suspension and sends cofiber
  sequences to long exact sequences.
\end{definition}

% SOURCE: aimpaper/main.tex:929-964 and all uses of distinguished triangles.
\begin{definition}[Cofiber and distinguished triangle data]
  \label{def:cofiber-triangle}
  \uses{def:stable-context}
  \notready
  For every map $f:X\to Y$, choose a cofiber $C(f)$ and a distinguished
  triangle
  \[
    X\xrightarrow fY\longrightarrow C(f)\longrightarrow\Sigma X.
  \]
  A morphism of triangles is three compatible maps, and a triangle equivalence
  is a morphism whose three components are isomorphisms.  Rotations and
  suspensions preserve distinguished triangles.  Every distinguished triangle
  whose first map is $f$ also carries a chosen triangle equivalence to this
  selected cofiber triangle; transporting along that equivalence is explicit
  data rather than definitional equality of its third object with $C(f)$.
\end{definition}

% SOURCE: aimpaper/main.tex:253-271, 648-664, and 929-971; exactness used repeatedly.
\begin{theorem}[Stable homotopy long exact sequence]
  \label{thm:stable-homotopy-long-exact}
  \uses{def:stable-homotopy-group,def:cofiber-triangle,def:mathlib-homological-functor,thm:mathlib-homology-sequence-exact-one,thm:mathlib-homology-sequence-exact-two,thm:mathlib-homology-sequence-exact-three}
  \notready
  Every chosen cofiber triangle
  $X\to Y\to C(f)\to\Sigma X$ induces the natural long exact sequence
  \[
    \cdots\to\pi_{n+1}C(f)\to\pi_nX\to\pi_nY
      \to\pi_nC(f)\to\pi_{n-1}X\to\cdots,
  \]
  with connecting maps compatible with rotation, suspension, and morphisms of
  triangles.
\end{theorem}

\begin{proof}
  Equip the representable functor $[S^0,-]$ with Mathlib's homological-functor
  structure, apply its three homology-sequence exactness lemmas to every
  shift, and splice the rotated triangles.  Naturality follows from a morphism
  of triangles.
\end{proof}

% SOURCE: aimpaper/main.tex:1755-1921, smash diagrams in May's lemma and the Mahowald trick.
\begin{definition}[Smash of cofiber triangles]
  \label{def:smash-triangle}
  \uses{def:cofiber-triangle}
  \notready
  Smashing a distinguished triangle with any spectrum gives a distinguished
  triangle, naturally in both variables.  For two cofiber triangles the
  induced $3\times3$ diagram contains the canonical pullback/pushout square
  and connecting maps used to compare the two ways of reaching the double
  cofiber.
\end{definition}

\section{Homotopy classes and secondary operations}

% SOURCE: aimpaper/main.tex:253-263, 283-306, and 912-997.
\begin{definition}[Adams filtration of elements and maps]
  \label{def:adams-filtration}
  \uses{def:admissible-spectrum,def:hf2-homology,def:stable-homotopy-group,def:adams-tower-filtration}
  \notready
  The primitive notions are predicates $\operatorname{AF}(\alpha)\ge s$ and
  $\operatorname{AF}(f)\ge k$: respectively, $\alpha$ lies in filtration
  $F^s\pi_nX$, and $f$ factors through the $k$th Adams-tower stage.  Their
  suprema take values in $\mathbb Z\cup\{+\infty\}$ for elements and
  $\mathbb N\cup\{+\infty\}$ for maps; in particular the zero element and zero
  map have infinite filtration.  The paper's notation
  $\operatorname{AF}(-)=k$ abbreviates that the finite supremum is attained
  and equal to $k$.  The interface records monotonicity under sums,
  composition, and the shift induced on detected classes.
\end{definition}

% SOURCE: aimpaper/main.tex:307-343, 408-459, and 461-648.
\begin{definition}[Detection relation]
  \label{def:adams-detection}
  \uses{def:strong-convergence-detection,def:adams-filtration}
  \notready
  For a permanent Adams class $x\in E_\infty^{s,t}(X)$ and a homotopy class
  $\alpha\in\pi_{t-s}X$, write $x\leadsto\alpha$ when $x$ is the associated
  graded image of $\alpha$ in filtration $s$.  Detection is a relation rather
  than a function: a homotopy element can have several representatives after
  adding higher-filtration terms.
\end{definition}

% SOURCE: aimpaper/main.tex:2387-2616, Toda-bracket arguments.
\begin{definition}[Toda bracket interface]
  \label{def:toda-bracket}
  \uses{def:cofiber-triangle}
  \notready
  Given composable stable maps $a,b,c$ with chosen nullhomotopies of $ba$ and
  $cb$, the triple Toda bracket $\langle c,b,a\rangle$ is a coset of homotopy
  classes with its standard left and right indeterminacy.  The interface
  includes suspension, naturality, containment under changed nullhomotopies,
  and Moss-style comparison with Massey products when an Adams convergence
  theorem supplies that comparison.
\end{definition}

% SOURCE: aimpaper/main.tex:2470-2567, Massey products on Adams pages; Moss Theorem 1.2.
\begin{definition}[Massey-product interface]
  \label{def:massey-product}
  \uses{def:multiplicative-ss,def:mathlib-f2}
  \notready
  On a differential graded page model for the multiplicative Adams spectral
  sequence, composable classes $a,b,c$ with $ab=0$ and $bc=0$ and chosen
  nullhomologies define a triple Massey-product coset
  $\langle a,b,c\rangle$.  The interface records its bidegree, left and right
  indeterminacy, naturality, change-of-choice formula, and a proposition that
  the indeterminacy is zero in a specified finite bidegree.
\end{definition}

% SOURCE: aimpaper/main.tex:2493-2567 and cited Toda/Moss identities; explicit load-bearing identities.
\begin{theorem}[Toda product and shuffle identities]
  \label{thm:toda-product-identities}
  \uses{def:toda-bracket,def:massey-product,def:adams-detection,def:external-result}
  \notready
  Explicitly located stable-homotopy results supply the Toda naturality,
  suspension, product, and shuffle containments used in Section 7, with their
  full indeterminacy.  In the particular degrees and chosen nullhomotopies of
  the paper they supply
  \[
    \eta^2\in\langle[h_0],\eta,[h_0]\rangle
  \]
  and, for every
  $b\in\langle\lambda^3\alpha_1,[h_0],\eta\rangle$,
  \[
    b[h_0]=\eta\lambda^3\eta\alpha_1.
  \]
  The record also carries the zero-indeterminacy checks needed for these
  equalities; a general containment is never silently strengthened to an
  equality.
\end{theorem}

\begin{proof}
  Instantiate the standard Toda composition and shuffle theorems with the
  paper's chosen nullhomotopies, calculate both indeterminacy subgroups, and
  use the typed vanishing evidence in the two relevant stems to show they are
  zero.  Transport the Hopf relation and signs to the characteristic-two
  convention.
\end{proof}

% SOURCE: aimpaper/main.tex:2387-2477; external Moss convergence theorem is required.
\begin{theorem}[Moss convergence adapter]
  \label{thm:moss-convergence-adapter}
  \uses{def:toda-bracket,def:massey-product,thm:toda-product-identities,def:adams-detection,def:external-result}
  \notready
  Given an explicit external result supplying the relevant Moss convergence
  theorem and hypotheses excluding crossings and indeterminacy, a Massey
  product in the Adams $E_r$-page detecting a permanent class yields a
  corresponding element of the stated Toda bracket in stable homotopy.
\end{theorem}

\begin{proof}
  Unpack the external theorem at the chosen spectra and representatives,
  discharge its convergence and no-crossing hypotheses from the supplied
  spectral-sequence data, and transport the detected class through the
  associated-graded isomorphism.  Keep the resulting Toda indeterminacy
  explicit rather than choosing a canonical element.
\end{proof}
